\documentclass[letterpaper, 12pt]{article}

%%%%%%%%%%%%%%%%%%%%%%%%%%%%%
% DEFINITIONS
% Change those informations
% If you need umlauts you have to escape them, e.g. for an ü you have to write \"u
\gdef\mytitle{Ausarbeitung}
\gdef\mythema{Pool 6a}

\gdef\mysubject{GGP-Matura}
\gdef\mycourse{5BHIT 2015/16}
\gdef\myauthor{Michael Weinberger}

\gdef\myversion{1.0}
\gdef\mybegin{17. Mai 2016}
\gdef\myfinish{20. Mai 2016}

\gdef\mygrade{}
\gdef\myteacher{Betreuer: Kraus}
%
%%%%%%%%%%%%%%%%%%%%%%%%%%%%%

\input special/preamble.tex

\let\tempsection\section
\renewcommand\section[1]{\vspace{-0.3cm}\tempsection{#1}\vspace{-0.3cm}}
\WithSuffix\newcommand\section*[1]{\tempsection*{#1}}

\let\tempsubsection\subsection
\renewcommand\subsection[1]{\vspace{0cm}\tempsubsection{#1}\vspace{0cm}}

\let\tempsubsubsection\subsubsection
\renewcommand\subsubsection[1]{\vspace{0cm}\tempsubsubsection{#1}\vspace{0cm}}

\linespread{0.94}

\lhead{\mysubject}
\chead{}
\rhead{\bfseries\mythema}
\lfoot{\mycourse}
\cfoot{\thepage}
% Creative Commons license BY
% http://creativecommons.org/licenses/?lang=de
\rfoot{\ccby\hspace{2mm}\myauthor}
\renewcommand{\headrulewidth}{0.4pt}
\renewcommand{\footrulewidth}{0.4pt}

\begin{document}
\parindent 0pt
\parskip 6pt

\pagenumbering{Roman} 
%!TEX root=../laborprotokoll.tex

\begin{titlepage}

	\begin{figure}[!h]
		\begin{flushright}
			\includegraphics[width=0.3\linewidth]{images/jdIT_tgm.png}
		\end{flushright}
	\end{figure}

	\vspace{2.5cm} 

	{\begin{center} \bfseries\huge
			\rule{17.5cm}{0.1mm}  
			\\[5mm]
			\mytitle\\[5mm]
			\mythema\\
			\rule{17.5cm}{0.1mm}  
	\end{center}}

	{\begin{flushright} \bfseries\Large
			\vspace{2cm}
			\mysubject\\
			\mycourse\\[10mm]
			\myauthor\\[10mm]
	\end{flushright}}

	{\begin{table}[!h] \bfseries\normalsize
		\begin{tabularx}{\textwidth}{lXr @{\hspace{0mm}}}
			&& Version \myversion\\
			\mygrade && Begonnen am \mybegin\\
			\myteacher && Beendet am \myfinish\\
		\end{tabularx}
	\end{table}}

\end{titlepage}


\clearpage
\thispagestyle{empty}
\tableofcontents

\newpage
\pagenumbering{arabic}
\pagestyle{fancy}

%\vspace{-0.5cm}

\section{Pool 6a - Wechselwirkungen von Kultur, Gesellschaft und Wirtschaft in der Geschichte \cite{buch}}

\subsection{Flucht, Vertreibung und Deportation: Migration und weltweite Kriege im 20. Jahrhundert}

Der Erste und Zweite Weltkrieg sowie der folgende globale 'Kalte Krieg' bedeuteten tiefe Einschnitte in die weltpolitische Ordnung und die weltwirschaftlichen Verhältnisse. Europa verlor in kurzer Zeit die über Jahrhunderte errungene Position als globales und politisches Zentrum, die Kolonialreiche waren aus politischen und finanziellen Gründen nicht mehr zu halten. Den 'Kalten Krieg' erlebte das 'Alte Europa' kaum mehr als eigenständig operierender weltpolitischer Akteur. Die weltweiten kriegerischen bzw. kriegsähnlichen Konflikte des 20. Jahrhunderts und deren politische Folgen führten zu einer enormen Zunahme der Zwangswanderungen in Gestalt von Deportation und Zwangsarbeit in den Kriegswirtschaften, Evakuierung und Flucht aus den Kampfzonen ebenso wie Massenausweisung und Vertreibung nach Kriegsende.

\subsubsection{Der Erste Weltkrieg als Motor des Zwangswanderungsgeschehens}

Der Erste Weltkrieg fürhte als 'totaler Krieg' zu einem rapiden Anwachsen der militärischen Kapazitäten der Gegner. Ein Kennzeichen der daraus resultierenden neuen Konfliktdynamik war, dass die militärischen Operationen zum Teil innerhalb weniger Tage und Wochen Millionen von Zivilisten in den Kampfzonen entwurzelten. Allein in den ersten drei Monaten nach dem Angriff Deutschlands flohen beispielsweise 1.4 Millionen Belgier, also ein Fünftel der 1914 knapp 7 Millionen Menschen umfassenden Gesamtbevölkerung des Landes, in die Niederlande, nach Frankreich oder Großbritannien. Weitere Hunderttausende verließen fluchtartig die Kampfzonen in Nord- und Nordostfrankreich, deren Bevölkerung noch ein Jahr nach Kriegsende mit rund 2 Millionen erst wieder rund 40 Prozent des Vorkriegsstandes erreichte. \\
Die Kriegssituation erleichterte bzw. ermöglichte eine staatliche Politik der Zwangsmigration gegenüber missliebigen Minderheiten. Erst der beschleunigte Ausbau der Interventions- und Ordnungskapazitäten der Staaten im Krieg bot die administrativen Instrumente, um Massenausweisungen oder Massenvertreibungen durchzuführen. Darüber hinaus förderte der Erste Weltkrieg die Verbreitung extremer Nationalismen - Fremdenfeindlichkeit wurde lanciert und die Tendenz zur Ausgrenzung von Minderheiten verstärkt. \\
Ein Instrument zum staatlichen Umgang mit 'feindlichen Ausländern' bildete die Internierung. Nicht weniger als 400.000 von ihnen wurden in den kriegsführenden europäischen Staaten 1914-1918 als 'Zivilgefangene' in Lagern festgehalten, Zehntausende darüber hinaus unter Zwang repatriiert.

\clearpage
\bibliographystyle{unsrt}
\bibliography{Pool6_Weinb_5BHIT}

\end{document}
